\documentclass[a4paper,11pt]{article}

\usepackage[utf8]{inputenc}

\usepackage{minted}

\usepackage{graphicx}
\usepackage{caption}
\usepackage{subcaption}
\usepackage{placeins}
\usepackage{pgfplots}
\pgfplotsset{compat=1.18} 

\begin{document}

\title{
    \textbf{Reduce Part 1}
}
\author{Jonathan Karlsson}
\date{Spring Term 2024}

\maketitle

\section*{Introduction}

The task was to create recursive functions.
\section*{Implementation}

The first function that was implemented was length It was implemented by simply calling the same function recursivly, and then increasing an argument by one each time while also removing an element, finally when the list is empty the value is returned.
\begin{minted}{elixir}
def length(list) do
  length(list,0)
end
def length([_|t],len) do
  length(t,len+1)
end
def length([],len) do
  len
end
\end{minted}
All of the functions were done recursivly and all also had a function for the empty list in the same way as the length, these always return either a value or a list. Depending on what the function does. 
The next fucntion was the even function, this function first calls itself with an extra list, then values are appended to this list every time the first value in the original list is an even number. When the list is empty the new\_list is returned.
\begin{minted}{elixir}
def even([h|t],new_list) do
  case rem(h,2) do
   0-> even(t,new_list++[h])
   _ -> even(t,new_list)
end
end
\end{minted}
The next function the inc simply copies the original list to a new list while adding the specified value to the copied element.
\begin{minted}{elixir}
def inc([h|t],val,new_list) do
  inc(t,val,new_list++[(h+val)])
end
\end{minted}
The next function, the sum was done by adding the current element to a sum which starts as 0 and then calling itself again with the rest of the list. Then return the result when the list is empty.
\begin{minted}{elixir}
def sum(list) do
  sum(list,0)
end
def sum([h|t],result)do
  sum(t,result+h)
end
def sum([],result) do
  result
end
\end{minted}
The dec was done in the same way as the sum but subtracting the value instead of adding it. 
The same is true for the mul function which multiplied the values. 
The odd is similar to the even but it checks for a 1 instead of a 0. The rem is the same as the previous per item calculation functions such as the 
mul and add but it uses the rem function with the value instead of doing a calculation. 
The prod function is very similar to the sum but it sets the initial value to nil, then checks if the value is nil or not if it is the element is used if it was not the element multiplied with the current result value is used as the next result. 
\begin{minted}{elixir}
def prod(list) do
  prod(list,nil)
end
def prod([h|t],result)do
  case result do
    nil -> prod(t,h)
    _ -> prod(t,result*h)
  end
end
def prod([],result) do
  result
end
\end{minted}
The final function is the div which checks if the remainder is 0 when using the rem function on the current element and the given value, if it is 0 the current element is added to the new list. 

\begin{minted}{elixir}
def div(list,val) do
  div(list,val,[])
end
def div([h|t],val,new_list) do
  case rem(h,val) do
    0 -> div(t,val,new_list++[h])
    _ -> div(t,val,new_list)
  end

  end

def div([],val,new_list) do
  new_list
end
\end{minted}
These functions are all very similar and can be much simpler by using other functions which will be done in part 2. 

\end{document}
